\documentclass[11pt,letterpaper]{article}
\usepackage[T1]{fontenc}
\usepackage{lmodern}
\usepackage[margin=1.05in]{geometry}
\usepackage{amsmath,amssymb,amsthm,booktabs,microtype,xurl}
\usepackage[hidelinks,pdftitle={Conditional I0-independence for a 451-state Turing machine},
  pdfauthor={James Chang}]{hyperref}

\newtheorem{theorem}{Theorem}[section]
\newtheorem{lemma}[theorem]{Lemma}
\theoremstyle{definition}
\newtheorem{definition}[theorem]{Definition}

\newcommand{\ZF}{\mathrm{ZF}}
\newcommand{\ZFC}{\mathrm{ZFC}}
\newcommand{\PA}{\mathrm{PA}}
\newcommand{\Con}{\operatorname{Con}}
\newcommand{\Prv}{\operatorname{Pr}}
\newcommand{\N}{\mathbb{N}}
\newcommand{\cardB}{\mathsf{BC}}
\newcommand{\setread}[1]{#1^{\omega}}
\newcommand{\file}[1]{\texttt{\nolinkurl{#1}}}
\setlength{\emergencystretch}{2em}

\title{Conditional \texorpdfstring{$I_0$}{I0}-independence for a\\451-state Turing machine}
\author{James Chang}
\date{}

\begin{document}
\maketitle

\begin{abstract}
We verify a binary Turing machine with 451 working states whose
blank-tape halting is conditionally independent of $\ZFC+I_0$.
The machine searches for a contradiction in a calculus for set
theory with a Berkeley cardinal. Its additional axiom expresses
elementarity by comparing an assignment with the identity in a
family of truth sets. We prove in Lean~4 that the exact table
implements the search, and that Peano arithmetic proves that
nonhalting implies consistency of a specified arithmetic
presentation of the searched theory. Let $B$ be $\ZF$ plus a
Berkeley cardinal and let $A$ be the Craig presentation of the
arithmetic consequences of $\ZFC+I_0$ used in the proof. If $B$
proves $\Con(A)$ in its natural numbers and $\ZFC+I_0$ is
consistent, then $\ZFC+I_0$ cannot prove nonhalting. Consistency
of $B$ gives actual nonhalting. With arithmetic $\Sigma_1$-soundness
of $\ZFC+I_0$ we also obtain unprovability of halting. Both
unprovability statements descend to $\ZF$ with the same $I_0$ sentence.
\end{abstract}

\section{Statement of the result}\label{sec:result}

Let $M$ be the machine in Section~\ref{sec:literal}. We use the language
of membership with equality. Let $\cardB$ assert the existence of a
Berkeley cardinal, and write
\[
  B=\ZF+\cardB,\qquad T=\ZFC+I_0,\qquad T^-=\ZF+I_0.
\]
The symbol $I_0$ in these theories denotes the first-order sentence
of Section~\ref{sec:large}. For an arithmetic sentence $\sigma$,
write $\setread{\sigma}$ for its interpretation in the natural numbers
of a model of set theory.

Let $A$ be the Craig presentation of the arithmetic theory
\[
  \operatorname{Ar}(T)=\{\,\sigma : T\vdash \setread{\sigma}\,\},
\]
where $\sigma$ ranges over closed arithmetic sentences.
Section~\ref{sec:incompleteness} specifies this presentation and the
consistency sentence $\Con(A)$. Write $H$ for the $\Sigma_1$ arithmetic
sentence that a finite run of $M$ from the blank tape terminates.
Its interpretation and its relation to the literal table are given
in Section~\ref{sec:arithmetic}.

\begin{theorem}\label{thm:main}
For the exact 451-state table the following hold.
\begin{enumerate}
\item If $B$ is consistent, then $M$ runs forever from the blank tape.
\item If $T$ is consistent and $B\vdash\setread{\Con(A)}$, then
      $T\nvdash\neg\setread{H}$ and $T^-\nvdash\neg\setread{H}$.
\item If $B$ is consistent and $T$ is arithmetically $\Sigma_1$-sound,
      then $T\nvdash\setread{H}$ and $T^-\nvdash\setread{H}$.
\end{enumerate}
Thus consistency of $B$, arithmetic $\Sigma_1$-soundness of $T$,
and $B\vdash\setread{\Con(A)}$ suffice for nonhalting and two-sided
independence in both theories.
\end{theorem}

Here arithmetic $\Sigma_1$-soundness means that if $T$ proves
$\setread{\sigma}$ for a closed arithmetic $\Sigma_1$ sentence
$\sigma$, then $\N\models\sigma$. In particular, it implies consistency
of $T$. The premise $B\vdash\setread{\Con(A)}$ is the form of relative
consistency used throughout the argument. It refers to an internal
proof in $B$ and to the same proof predicate used in the application
of G\"odel's second incompleteness theorem.

The 451-state construction was obtained by compressing the register
program and decoder of a larger Berkeley-based proof search. The
verification below works directly with the 451-state table. Its two
directions use different forms of correctness. A halting computation
can be interpreted in a model of $B$ using the soundness of the
command calculus. For the nonhalting unprovability result, every
model of $\PA$ must admit the converse construction on its
internally finite proof codes.

\section{The literal machine}\label{sec:literal}

The tape is two-way infinite with alphabet $\{0,1\}$. State $0$ is
initial and the tape is entirely zero. There are 451 working states,
numbered $0,\ldots,450$, and a separate terminal label $451$.
For a state $q$ and symbol $b$, a row $(q,b,w,D,q')$ writes $w$,
moves one cell in direction $D\in\{L,R\}$, and enters $q'$.
The complete 902-row table is \file{data/zf-berkeley-451.tm} in the
companion source package. Its SHA-256 digest is
\begin{quote}\small
\nolinkurl{034c26fb655f32123c516bfab93be38d5e055415058eb29e856c4c9a80e74d64}.
\end{quote}
The first row is $(0,0,1,R,1)$. The unique row entering the terminal
label is $(156,1,1,R,451)$, and the terminal label has no outgoing
rows. These properties are checked on the packed literal table.

The initialization has two parts. After 39,753 literal steps the
head reaches decoder state $7$. At this point the tape has a left
marker, an encoded instruction pointer, a gap of zeros and a unary
register file. The initial frame has 192 fixed blocks and nine live
registers, all nine initially zero. A block $0\,1^{v+1}$ holds $v$.
The distance from the instruction-pointer start to the first
register is called the slot.

The next 126 interpreter instructions bring the slot from $19$ to
$107$ and enter the region used by the permanent interpreter
invariant. These instructions are checked through the interpreter
contracts. The loader, by comparison, is checked by concatenating
finite literal runs. At a decoder entry a word and a padding of zeros
specify the address. Explicit trailing zeros of the word can be shifted
into the padding at fixed slot, since the physical tape is unaltered.

The decoder proof gives a rank for each active control and two-bit
left context. It verifies the local effect on both read symbols,
including contractions of up to three tape steps. Each recursive
case lowers the rank, and at most 94 address bits are needed at
the root. A separate condition on zero runs and jump markers verifies
the return paths. The register contracts compute the effects of
increment and of both decrement branches on arbitrary unary blocks.
The space invariant combines these results so that decoding, register
access and continuation preserve the permitted frame at the next
abstract boundary. This is the part of the proof that connects the
finite transition checks to an unbounded run.

\section{The proof search}\label{sec:search}

\subsection{Formula and command codes}

Let $p(a,b)=\binom{a+b+1}{2}+b$ be Cantor pairing. Implication,
membership, universal quantification and equality have constructors
\[
\begin{array}{ll}
  \mathsf{imp}(a,b)=4p(a,b), & \mathsf{mem}(i,j)=4p(i,j)+1,\\
  \mathsf{all}(i,a)=4p(i,a)+2, & \mathsf{eq}(i,j)=4p(i,j)+3.
\end{array}
\]
Here $i,j$ are variable labels. Code $0$ denotes truth. Code $1$
denotes $v_0\in v_0$, the contradiction used by the machine. The
calculus includes irreflexivity of membership, so this atom can
serve in the coding of negation.

At a command boundary the proof store consists of the current formula,
a stack coded by nested pairs, a numeric operand and a cached formula.
A successful \textsc{mp} tests that the cached formula is the implication
from the current formula to the candidate consequence.
\textsc{drop} exposes the cached formula and saves the old current
formula on the stack. \textsc{nip} pops the stack into the cache,
except that it terminates when the cache already contains code $1$.
\textsc{extra} saves the current formula and places the Foundation
sentence and the compressed Berkeley sentence above it. The
\textsc{ax} routine builds the axiom bundle and the guarded
logical operations. The implementation follows the numerical
proof-search design of the \texttt{m254}/\texttt{tg296} development
\cite{ketchersid}.

A command word is stored as a nested pair code. A fresh candidate
is read directly. Between its commands the queued tail is stored
with a one-unit delimiter. When the queue is empty the machine
loads the next natural-number candidate. Otherwise it subtracts
one before unpacking the next command. The remainder of an
unpacked code is at most that code, so a nonempty queue strictly
decreases. In a nonhalting run each candidate is therefore reached.
Packing a finite command word gives a candidate that executes the
word in order. The bootstrapped interpreter and the command-boundary
invariant connect this enumeration to the literal run.

\subsection{The calculus searched by the machine}

Write $P$ for the compressed Berkeley sentence emitted by
\textsc{extra}. The nonlogical base used by the calculus includes
Extensionality and Foundation, together with a container principle
and a collection scheme. In the container principle every $a$
belongs to a transitive set $D$ such that for each $x\in D$ there
is a $w\in D$ containing every subset of $x$. In the collection scheme,
given a formula $R(x,y,\bar z)$ and a set of inputs, an image set
contains a witness for each input that has one. Each element of the
image is in the global range of $R$.

In a model of $\ZF$ the container can be taken at a sufficiently
high limit rank. Collection follows by bounding the least ranks
of witnesses for the inputs with witnesses, and cutting out the
required elements from a common rank bound. For the reverse
interpretation the collection scheme yields Separation and
Replacement. The container bounds unions and power sets, and the
ordinal members of a container of the empty set give an inductive
set. With the other base axioms these principles suffice for
$\ZF$. The model-theoretic verification applies also to arbitrary
first-order models of the base.

The command soundness invariant asserts that the top, cache, and
each stack entry are derivable. The operand has a purely numeric
role. For completeness, the detachment and axiom-building contracts
give command words for derivations, using a single-formula push
contract over a preserved stack. For example the closed formulas
produced by \textsc{extra} can each be isolated by stack and cache
operations. A proof of code $1$ can place that formula in the cache
and invoke \textsc{nip}. The resulting standard-natural theorem is
\begin{equation}\label{eq:search}
  M\text{ halts on the blank tape}
  \quad\Longleftrightarrow\quad
  1\text{ is derivable in the native calculus}.
\end{equation}
The Lean declaration for \eqref{eq:search} is
\file{I0451.Runtime.literal_halts_iff_derives}.

\section{The large-cardinal sentences}\label{sec:large}

\subsection{The Berkeley sentence emitted by the machine}

A Berkeley cardinal $\delta$ is an uncountable cardinal such that
for every transitive set $D$ with $\delta\in D$ and every $\eta<\delta$,
there is an elementary $j:D\to D$ with
$\eta<\operatorname{crit}(j)<\delta$ \cite{bkw}. The emitted formula
uses a different witness. It asks for a set $a$ such that every
transitive $D$ containing $a$ admits an elementary self-map moving
some member of $a$, with elementarity expressed in the following
compressed form.

An assignment on $D$ is a map $D\to D$, using the members of $D$
as labels. Pair graphs give a set $U$ of such assignments. For
$i,k\in D$, begin with the truth set
\[
 E_{ik}=\{s\in U:s(i)\in s(k)\}.
\]
For $R,S\subseteq U$ and $i\in D$, the closure operation takes
\begin{equation}\label{eq:cyl}
  \{s\in U:\exists t\in U\
       [\,s\equiv_i t\ \land\ (t\in R\ \longrightarrow\ t\in S)\,]\},
\end{equation}
where $s\equiv_i t$ means agreement outside label $i$. Representatives
in the literal formula are compared by their traces on $U$.
There is a closed family containing the atomic truth sets and
closed under \eqref{eq:cyl}. If $e$ is the identity assignment and
$j$ is the proposed map, the literal condition requires that each
member of the family whose trace contains $e$ also contain $j$.

\begin{lemma}\label{lem:elementarity}
For an infinite transitive set $D$, the truth-set condition implies
that $j$ is elementary. An elementary self-map of $D$ satisfies the
condition with a suitable family of traces.
\end{lemma}

\begin{proof}
Irreflexivity makes the diagonal atom the empty truth set. Applying
\eqref{eq:cyl} with empty antecedent gives $U$. For two truth sets
of finite support, choose $i$ outside the two supports and obtain
ordinary implication. Implication to the empty truth set gives
complement. Applying \eqref{eq:cyl} with antecedent $U$ gives
existential quantification at $i$. Equality is expressed by agreement
of membership extensions in $D$, using its transitivity and
Extensionality. Induction now represents every internally finite
formula.

For a parameter tuple, rename bound labels away from the finitely
many parameter values. Use the parameter values themselves as
the free labels, with a shared label for repeated values. Evaluating the
resulting formula at $e$ gives its truth at the original tuple.
Evaluating at $j$ gives truth at its image. Applying the comparison
to a formula and its negation gives both implications. Conversely,
take the traces represented by internally finite formulas.
The closure operations preserve this family and elementarity
gives the comparison at $e$ and $j$.
\end{proof}

The formalization constructs the least closed trace family by
Separation in the power set of $U$. It identifies that family with
finite stages of the closure operations and uses coded satisfaction
to implement the renaming argument inside the model. In particular,
Lemma~\ref{lem:elementarity} is used for internally finite codes
rather than merely for each externally standard formula.

\begin{lemma}\label{lem:berkeley}
Over $\ZF$, the emitted sentence $P$ is equivalent to $\cardB$.
\end{lemma}

\begin{proof}
If $\delta$ is Berkeley, take $a=\delta$. The defining property
supplies the required moving map on every transitive set containing
$\delta$, and Lemma~\ref{lem:elementarity} supplies its comparison
family.

Conversely, take a witness $a$ and put $\gamma=\operatorname{rank}(a)$.
Given a transitive $M$ with $\gamma+1\in M$, choose a sufficiently
high limit rank $V_\rho$. Include $M$ and $V_\rho$ as a pair in a
larger transitive set $N$ in which the pair is definable. One can
do this by adding a set of pairs with fixed first coordinate to a
limit rank, making the added set the unique membership-maximal
element. See \cite[Lemma~3.1]{bkw}. The map supplied by $P$ on $N$
moves a member of $a$ and fixes $M$ and $V_\rho$ setwise.
Relativization gives elementary restrictions. If all ordinals
below $\gamma$ were fixed, induction on ranks in $V_\rho$ would
fix $V_\gamma$ pointwise, contradicting the moved member of $a$.
The least moved ordinal is therefore below $\gamma+1$ and is
the critical point of the restriction to $M$. This makes
$\gamma+1$ proto-Berkeley. The least proto-Berkeley ordinal is
an uncountable cardinal with the required cofinality of critical
points \cite[Theorem~3.3 and Remark~3.5]{bkw}. The formalization
checks these leastness arguments in the coded set language.
\end{proof}

Together with the base axioms, Lemma~\ref{lem:berkeley} gives
both uses of $B$ needed below. Its models validate the native
calculus, and the native calculus proves the interpreted axioms
of $B$.

\subsection{The first-order presentation of \texorpdfstring{$I_0$}{I0}}

The target axiom is expressed using the proper-ultrapower
presentation of \cite[Definitions~2.5 and~2.7 and Proposition~3.5]{schlutzenberg}.
It uses an even ordinal $\lambda$, the set $X=V_{\lambda+1}$,
and a set ultrafilter $\mu$ on the subsets of $X$ in $L(X)$.
The ultrapower uses functions in $L(X)$ with domain $X$.
For each function $f$, the collapse of its quotient value is
specified by a set of representatives containing $f$ and covering
the predecessors modulo $\mu$. A transitive set is the range of
the collapse map, with equality and membership on it agreeing
with the corresponding relations modulo $\mu$.
The sentence requires compatibility of the collapses and coverage
of $L(X)$.

The induced constant-function map $j$ fixes $X$ and $\lambda$.
Its critical point is below $\lambda$, and $j\upharpoonright V_\lambda$
belongs to $L(X)$. It takes each relative constructible stage to
the stage at the image ordinal and preserves internally coded
satisfaction between those two set stages. The relative hierarchy
starts at the transitive closure of $\{X\}$, takes definable subsets
at successors, and takes unions at limits.

For a proper $I_0$ witness, the derived ultrafilter gives the
quotient collapse and the elementary constant-function map.
Collection supplies the representatives for the local collapse
of a given value. Since the stage recursion is uniform in $X$,
elementarity gives the two stage clauses. Conversely, the local
collapses give an isomorphism of the quotient with $L(X)$.
The stages are transitive and exhaust that class. Their satisfaction
clauses give bounded elementarity, and the ordinal images are
cofinal. Mapping stages to their image stages therefore gives
membership cofinality. Cofinal bounded elementarity yields full
elementarity. This gives the mathematical identification of the
coded sentence with the proper-witness presentation. The Lean
formalization encodes the sentence and proves the result for
the theories thereby specified.

\section{Arithmetizing the literal computation}\label{sec:arithmetic}

In a model of $\PA$, represent a configuration by $(q,\ell,r)$.
The two tape stacks $\ell$ and $r$ are binary integers with the nearest
cell in the low bit. The right stack starts at the head. The table is
one integer whose base-$2048$ digit at $2q+b$ encodes the row for
state $q$ and read symbol $b$. A row with target $q'$, direction
bit $d$ and write symbol $w$ has digit $1+4q'+2d+w$.
The bit $d=1$ denotes a right move. With integer division, the tape
updates for the two directions are
\[
\begin{array}{ll}
 R: & (q,\ell,r)\longmapsto (q',2\ell+w,\lfloor r/2\rfloor),\\
 L: & (q,\ell,r)\longmapsto
 (q',\lfloor\ell/2\rfloor,4\lfloor r/2\rfloor+2w+(\ell\bmod 2)).
\end{array}
\]
A zero digit denotes the absence of a move, in which case the
arithmetized transition leaves the configuration fixed. The formula
$H$ asserts that the primitive recursive iteration from $(0,0,0)$
reaches a configuration with a zero digit. It is $\Sigma_1$.
The standard-model equivalence is
\[
   \N\models H \quad\Longleftrightarrow\quad
   M\text{ halts on the blank tape}.
\]
The interpretation $\setread H$ reads this assertion in the internal
$\omega$ of any model of $\ZF$.

The direction used for incompleteness needs witnesses that a
nonstandard model of $\PA$ regards as finite. The native arithmetic
calculus uses a proof book with coded axiom, generalization and
detachment lines. A proof book for code $1$ gives a coded command
word by the stack contracts. The coded queue decreases between
candidate loads, so an internally finite word is reached and
executed through arithmetic induction. Register updates are
transported through arithmetic histories of the literal tape by
the arithmetized decoder and register contracts. This proves
that an internally finite native refutation produces an
internally finite literal halting history.

To apply this argument to ordinary arithmetical proof predicates,
let $S$ be the fixed arithmetic presentation obtained from the
native calculus. A set-language presentation enumerates the
sentences whose translations into the native syntax have native
proof books. An arithmetic presentation enumerates sentences
whose readings in the set language belong to that enumeration.
Both recognizers have explicit $\Sigma_1$ definitions. Padding
an axiom with a coded witnessing computation gives a
$\Delta_1$ axiom predicate. There are two relevant facts.
\begin{enumerate}
\item If $B\vdash\setread{\sigma}$, then $S\vdash\sigma$.
      The native set presentation proves the axioms of $B$.
      A coded sequent-to-Hilbert translation compiles its proof
      of the reading of $\sigma$ to a native derivation.
\item In any model of $\PA$, an $S$-refutation gives a native
      proof book for code $1$. The coded arithmetic-to-set
      interpretation transports the axiom cases and the logical
      rules. The translated proof produces a contradiction in the
      native set presentation, which the native compiler transports
      to a book of the required form.
\end{enumerate}
The first fact is a metatheoretic statement about proofs. In
particular $S$ contains the theorems of $\PA$. The second fact
holds for internally finite proofs. Combining it with the
arithmetized run above and applying completeness for arithmetic
gives the object-theory statement
\begin{equation}\label{eq:internal}
  \PA\vdash \neg H\longrightarrow\Con(S).
\end{equation}
This is \file{I0451.literal_nonhalting_implies_native_consistency}.
The proof of \eqref{eq:internal} uses the actual literal run
predicate. The distinction between \eqref{eq:search} and
\eqref{eq:internal} is essential because standard computations alone
do not establish an arithmetic implication in every model.

\section{The incompleteness argument}\label{sec:incompleteness}

The theory $\operatorname{Ar}(T)$ extends $\PA$. To enumerate its
theorems, use the finite-parameter versions of the $\ZF$ schemata
and a coded reading of arithmetic sentences in $\omega$.
The search for a proof of a reading is recursively enumerable.
The equivalence of these readings with the interpretation
$\setread{\sigma}$ gives exactly $\operatorname{Ar}(T)$.
Foundation's Craig construction uses a chosen standard-natural
recognizer for this enumeration and pads its axioms with
bounded witnesses. This yields the specific $\Delta_1$
presentation $A$ used here. Its standard theorem set agrees with
$\operatorname{Ar}(T)$, so consistency of $T$ implies consistency
of $A$. We write $\Con(A)$ for $\neg\Prv_A(\ulcorner\bot\urcorner)$
using this proof predicate, and likewise for $S$.

\begin{lemma}\label{lem:transfer}
If $B\vdash\setread{\Con(A)}$, then
\[
        \PA\vdash \Con(S)\longrightarrow\Con(A).
\]
\end{lemma}

\begin{proof}
The first proof transport of Section~\ref{sec:arithmetic} gives
$S\vdash\Con(A)$. Since $\neg\Con(A)$ is $\Sigma_1$,
formalized $\Sigma_1$ completeness gives, in $I\Sigma_1$,
\[
  \neg\Con(A)\longrightarrow\Prv_S(\ulcorner\neg\Con(A)\urcorner).
\]
The fixed $S$-proof of $\Con(A)$ turns a coded $S$-proof of
its negation into a coded $S$-refutation. This gives
$\neg\Con(A)\to\neg\Con(S)$ within $I\Sigma_1$. Contraposition
and extension to $\PA$ prove the lemma.
\end{proof}

\begin{proof}[Proof of Theorem~\ref{thm:main}]
If $B$ is consistent, first-order completeness gives a model
of $B$. It validates the axioms checked by the native calculus,
by the model argument in Section~\ref{sec:search} and
Lemma~\ref{lem:berkeley}. Code $1$ is false in that model.
By \eqref{eq:search} the machine cannot halt.

For the second assertion, compose \eqref{eq:internal} with
Lemma~\ref{lem:transfer}. The result is
\[
         \PA\vdash \neg H\longrightarrow\Con(A).
\]
A proof of $\neg\setread H$ in $T$ would give a proof of
$\neg H$ in $A$. Hence $A$ would prove $\Con(A)$, contrary to
G\"odel's second incompleteness theorem for this presentation
\cite{goedel}. For the third assertion, $H$ is false in $\N$ by the
first assertion and the exact arithmetic interpretation. Arithmetic
$\Sigma_1$-soundness excludes a $T$-proof of $\setread H$.
Both unprovability statements descend by inclusion $T^-\subseteq T$.
\end{proof}

\paragraph{Scope of the consistency premise.}
Goldberg attributes to Woodin the result that $\ZF$ plus a
Berkeley cardinal proves the consistency of $\ZFC+I_0$
\cite[slide~15]{goldberg}. The premise used here is a proof in
$B$ for the particular Craig predicate just specified.
Discharging it from the conventional set-theory consistency
statement would require an implication provable in $B$ between
the two internal consistency sentences. The equality of the
standard theorem sets gives external equiconsistency. That
calculation leaves the internal implication open. Theorem~\ref{thm:main}
states the premise for the presentation used in its proof.

\section{Formal verification}\label{sec:verification}

The main Lean declaration is \file{I0451.exact451_independence}.
The unprovability of nonhalting with the weaker consistency
assumptions is \file{I0451.i0_nonhalting_unprovable}.
The standalone statement \file{I0451.Standalone.exact451Claim}
uses only Lean's standard library in its Challenge.
It gives the transition table and the membership-model statement.
Comparator checks the claim together with the four formula values
supplied by the Solution, for $\cardB$, $I_0$, the consistency
premise and the arithmetic halting sentence. The alignment proof
uses first-order completeness to connect validity with provability
and verifies the interpretation of $\Sigma_1$ soundness.

The proof reuses Foundation's arithmetized syntax, proof predicates
and incompleteness arguments \cite{foundation}, and the arithmetic
interpretation and proof-calculus libraries shared with the
207-state development \cite{zf207}. The checked table and decoder
data are generated from the machine representation. The simulation,
logical interpretation and proof transports are reusable Lean
arguments. The local build and axiom audit passed with Lean~4.35.0.
The audited independence declarations use only \file{propext},
\file{Classical.choice} and \file{Quot.sound}. Lean's kernel and
NanoDa accepted the Comparator export, and the bundled Con-Ron
accepted the same export in verified mode. The source package
records the reproduction commands and includes the literal and named
transition tables.

\paragraph{Acknowledgments.}
The Lean development and preparation of this manuscript used Codex.
The source package documents the construction and library provenance.

\begin{thebibliography}{9}

\bibitem{bkw}
J.~Bagaria, P.~Koellner, and W.~H.~Woodin.
\newblock Large cardinals beyond choice.
\newblock \emph{Bulletin of Symbolic Logic}, 25(3):283--318, 2019.
\newblock \url{https://doi.org/10.1017/bsl.2019.28}.

\bibitem{zf207}
J.~Chang.
\newblock The ZF207 proof package.
\newblock \url{https://github.com/jamesyc/BusyBeaver/tree/9437ef874e1c76a0ec45330554e895157d56c604/ZF207}.

\bibitem{foundation}
Formalized Formal Logic.
\newblock \emph{Foundation}, Lean library, revision \file{b2d8f55}.
\newblock \url{https://github.com/FormalizedFormalLogic/Foundation}.

\bibitem{goedel}
K.~G\"odel.
\newblock \"Uber formal unentscheidbare S\"atze der Principia Mathematica
und verwandter Systeme I.
\newblock \emph{Monatshefte f\"ur Mathematik und Physik}, 38:173--198, 1931.

\bibitem{goldberg}
G.~Goldberg.
\newblock Introduction to choiceless large cardinals.
\newblock Lecture slides, 2022.
\newblock \url{https://advances.assafrinot.com/wp-content/uploads/2022/07/goldberg1.pdf}.

\bibitem{ketchersid}
J.~Ketchersid.
\newblock \texttt{m254} and \texttt{tg296}, \emph{busy-beaver}, revision
\file{cde3649}.
\newblock \url{https://github.com/jeff87654/busy-beaver/tree/cde3649e0f1f3ae57454be7a5bed8bf19f8115d7}.

\bibitem{schlutzenberg}
F.~Schlutzenberg.
\newblock On the consistency of ZF with an elementary embedding from
$V_{\lambda+2}$ into $V_{\lambda+2}$.
\newblock \emph{Journal of Mathematical Logic}, 2024, online first.
\newblock \url{https://arxiv.org/abs/2006.01077v4}.

\end{thebibliography}
\end{document}
